% ---------------------------------------------------------------
% Subjective evaluation subsection. Drafted 2026-09-22, rewritten
% 2026-09-23 (by request: the official name is "subjective evaluation";
% declarative statements only, no narrative). Numbers are from
% ratings_7.csv through subjective_spider.py --drop-constant (17 raters
% kept), NOT remembered:
%   demographics, per-system means      subjective_spider.py stdout
%   paired differences and p-values     results/subjective_paired.csv
%   the table                           paper/tables/listening_mos.tex
%   the spider chart (optional figure)  figures/subjective_spider.pdf
% System names as in experiments.tex. The cascade and the causal-
% attention arm were not rated.
%
% FILTERING, stated in the text: 18 rows from a 14-15 Sept pilot used an
% earlier page whose system codes A-G are not recoverable, dropped; four
% raters gave one identical score to every clip on every axis (three
% aged 10, none with musical training), dropped. Their removal changes
% no p-value: a constant rater contributes only ties, which the
% signed-rank test discards.
%
% PLACEHOLDERS [X] must not be guessed:
%   - the exact wording of the five questions on the rating page (the
%     page is not in the repo; the axis names are the CSV columns
%     consistency / structure / fit / musicality / creativity);
%   - clip content and length in bars: the five-bar prompt plus the
%     continuation, rendered with one piano/strings timbre pair for every
%     system (render_piano_copies.py);
%   - whether the prompt was played before the continuation; whether
%     the clip order was randomised per song;
%   - how raters were recruited.
% Preamble: amsmath only; no \paragraph headings.
% ---------------------------------------------------------------
\subsection{Subjective evaluation}
\label{sec:subjective}

\textbf{Material.} Five held-out POP909 songs, cropped as in
Sec.~\ref{sec:protocol}. For each song, six clips: the continuation of
each co-generation system of Sec.~\ref{sec:systems} and the
ground-truth continuation. All clips of a song are rendered with the
same piano timbre for melody and the same string timbre for chord [X:
clip length; whether the prompt is included].

\textbf{Procedure.} A rater is assigned a song and rates its six clips
[X: order randomised]. Each clip is rated on five axes on a scale of
$1$ to $5$: \emph{consistency}, coherence of the continuation with the
prompt; \emph{structure}, phrase organisation and development over the
clip; \emph{melody--chord fit}, agreement between the two streams;
\emph{musicality}; \emph{creativity} [X: exact question wording]. A
rater may rate further songs. Ratings from a pilot on an earlier
version of the page are excluded.

\textbf{Participants.} 21 raters produced 438 ratings [X:
recruitment]. Four raters assigned one identical score to every clip
on every axis and are excluded; their exclusion changes no
significance result reported below, because a constant rater
contributes only tied pairs. The remaining 17 raters (10 male, 7
female; age 21--55, median 27) report some musical training (11), none
(4), or amateur musicianship (2). They contribute 372 ratings, 62 per
system, 6--30 per rater (median 30); each song is rated by 11--15
raters.

\textbf{Statistics.} Each rater rates every system on a song, so a
rater--song pair is the unit of comparison. For each system and each
axis, and for the mean over axes, the difference between the rating of
\sys and the rating of the system on the same rater--song pair is
tested with a two-sided Wilcoxon signed-rank test over the 62 pairs;
tied pairs are discarded. The bootstrap intervals reported with the
means resample raters, since one rater's ratings are not independent.

\begin{table}[t]
\caption{Listening test: mean rating (1--5) per system and axis over 17 raters and five prompts; each rater heard every system on a prompt. $^{\mathrm{+}}$/$^{\mathrm{-}}$: rated significantly above / below \emph{Duet (ours)} (two-sided Wilcoxon signed-rank over rater--prompt pairs, $p<0.05$). Bold: best system per axis, ground truth aside.}
\label{tab:listening}
\begin{center}
\footnotesize
\setlength{\tabcolsep}{4pt}
\begin{tabular}{|l|c|c|c|c|c|c|}
\hline
\textbf{System} & \textbf{Consistency} & \textbf{Structure} & \textbf{Fit} & \textbf{Musicality} & \textbf{Creativity} & \textbf{Overall} \\
\hline
Duet (ours) & \textbf{3.92} & \textbf{3.85} & \textbf{3.94} & 3.65 & \textbf{3.71} & \textbf{3.81} \\
Single-stream (scratch) & 3.08$^{\mathrm{-}}$ & 3.05$^{\mathrm{-}}$ & 2.98$^{\mathrm{-}}$ & 2.82$^{\mathrm{-}}$ & 3.00$^{\mathrm{-}}$ & 2.99$^{\mathrm{-}}$ \\
Whole-Song Gen & 3.26$^{\mathrm{-}}$ & 3.39$^{\mathrm{-}}$ & 3.45$^{\mathrm{-}}$ & 3.26$^{\mathrm{-}}$ & 3.40$^{\mathrm{-}}$ & 3.35$^{\mathrm{-}}$ \\
Anticipatory Music Transf. & 3.61$^{\mathrm{-}}$ & 3.60 & 3.79 & 3.60 & 3.66 & 3.65 \\
Single-stream (finetuned) & 3.81 & 3.66 & 3.73 & \textbf{3.77} & 3.58 & 3.71 \\
Ground truth & 4.18$^{\mathrm{+}}$ & 4.26$^{\mathrm{+}}$ & 4.39$^{\mathrm{+}}$ & 4.26$^{\mathrm{+}}$ & 4.11$^{\mathrm{+}}$ & 4.24$^{\mathrm{+}}$ \\
\hline
\end{tabular}
\end{center}
\end{table}


\textbf{Results.} Table~\ref{tab:listening} reports the means. \sys
obtains the highest mean among the systems on consistency ($3.92$),
structure ($3.85$), melody--chord fit ($3.94$) and creativity ($3.71$),
and overall ($3.81$); the finetuned single-stream model obtains the
highest mean on musicality ($3.77$ against $3.65$). Paired
differences, \sys minus system, are as follows.
\begin{itemize}
  \item Single-stream, scratch ($2.99$ overall): $+0.83$ overall,
    $p < 0.001$; positive and significant on every axis
    ($+0.71$ to $+0.95$, $p < 0.001$); \sys rated higher in 47 of 62
    pairs and lower in 10.
  \item Whole-Song-Gen ($3.35$): $+0.46$ overall, $p < 0.001$;
    positive and significant on every axis, from $+0.31$ on creativity
    ($p = 0.036$) to $+0.66$ on consistency ($p < 0.001$).
  \item Anticipatory Music Transformer ($3.65$): $+0.16$ overall,
    $p = 0.12$; significant on consistency only ($+0.31$, $p = 0.031$);
    $+0.05$ on musicality and on creativity ($p > 0.6$).
  \item Single-stream, finetuned ($3.71$): $+0.10$ overall,
    $p = 0.33$; no axis is significant; the largest differences are
    $+0.21$ on fit ($p = 0.10$) and $-0.13$ on musicality
    ($p = 0.36$).
  \item Ground truth ($4.24$): $-0.43$ overall, $p < 0.001$; negative
    and significant on every axis, from $-0.26$ on consistency
    ($p = 0.037$) to $-0.61$ on musicality ($p < 0.001$).
\end{itemize}
The subjective ordering, ground truth, \sys, finetuned single-stream,
Anticipatory Music Transformer, Whole-Song-Gen, scratch single-stream,
agrees with the objective ordering of Table~\ref{tab:e1}. The
difference between \sys and the two external systems is largest on
consistency and smallest on musicality and creativity. The difference
between \sys and the finetuned single-stream model, which shares its
backbone, training songs and tokenisation, is not significant on any
axis.
